\Needspace{16\baselineskip}
\fancyhead[L]{\small VESTA.}
\begin{center}
VESTA.

\par\vspace{4pt}\hrule\vspace{4pt}

\fontsize{9}{11}\selectfont Göttingische gelehrte Anzeigen. Stück 55. Seite 537..540. 8. April 1809.

\par\vspace{4pt}\hrule\vspace{4pt}
\end{center}

In einigen früheren Ausgaben unserer Gel. Anz. haben wir die ersten Beobachtungen der Vesta mitgeteilt, die von Juni bis September des vergangenen Jahres an der hiesigen Sternwarte gemacht wurden, ebenso wie die Beobachtungen von Professor \textsc{Bode} in Berlin (1808, S. 1793) [Nov. 10. S. 308 dieses Bandes]. Vor kurzem hat Professor \textsc{Gauss} von Herrn \SourceLeaf{315}\textsc{Bouvard} auch die Beobachtungen dieses Planeten erhalten, die an der kaiserlichen Sternwarte in Paris gemacht wurden. Wir glauben, den Astronomen einen Dienst zu erweisen, wenn wir sie frühzeitig bekannt machen.

\begin{center}
Beobachtungen der Vesta in Paris:
\end{center}


\begin{center}\fontsize{10}{12}\selectfont\setlength{\tabcolsep}{5pt}\renewcommand{\arraystretch}{1.15}\begin{tabular}{lrrr}
1808 & \shortstack{Mittlere Zeit\\in Paris} & \shortstack{Scheinbare\\Rektaszension} & \shortstack{Scheinbare südliche\\Deklination}\\
\hline
Aug. 26 & $13^{\text{u}}\;16^{\text{m}}\;23^{\text{s}}5$ & $354^\circ\;11'\;28''80$ & $14^\circ\;1'\;20''3$\\
28 & $13\;6\;56.1$ & $353\;47\;30.30$ & $14\;17\;43.4$\\
Sept. 16 & $11\;35\;23.2$ & $349\;34\;9.45$ & $16\;34\;8.7$\\
19 & $11\;20\;57.3$ & $348\;54\;30.00$ & $16\;50\;16.0$\\
22 & $11\;6\;37.0$ & $348\;16\;15.00$ & $17\;4\;24.0$\\
24 & $10\;57\;7.7$ & $347\;51\;49.80$ & $17\;12\;48.5$\\
25 & $10\;52\;24.5$ & $347\;39\;55.35$ & $17\;16\;36.0$\\
26 & $10\;47\;42.1$ & $347\;28\;18.75$ & $17\;20\;10.0$\\
Oct.  4 & $10\;10\;48.6$ & $346\;6\;32.55$ & $17\;39\;50.5$\\
5 & $10\;6\;18.0$ & $345\;57\;50.25$ & $17\;41\;6.7$\\
6 & $10\;1\;48.7$ & $345\;49\;27.75$ & $17\;42\;20.0$\\
13 & $9\;31\;7.4$ & $345\;1\;51.75$ & $17\;43\;36.3$\\
\end{tabular}\end{center}


{\rightskip=0pt plus 1em \par\indent Aus diesen und den übrigen vorhandenen Beobachtungen hat Professor \textsc{Gauss} die Opposition der Vesta bestimmt, die \emph{erste}, die bisher beobachtet wurde.\par}

\Needspace{7\baselineskip}
\begin{center}Opposition der Vesta:

1808 September 8. $7^{\text{u}}\;52^{\text{m}}\;57^{\text{s}}$ Mittl. Zeit in Göttingen.
\end{center}
\begin{center}\fontsize{10}{12}\selectfont\setlength{\tabcolsep}{5pt}\renewcommand{\arraystretch}{1.15}\begin{tabular}{lr}

Wahre Länge .\;.\;.\;.\;.\;.\;.\;.\;.\; & $345^\circ\;53'\;38''5$\\
Wahre geozentrische Breite .\;.\;.\;.\; & $11\;0\;25.8$ südl.\\
\end{tabular}\end{center}


\par\indent An der hiesigen Sternwarte wurden die Meridianbeobachtungen der Vesta bis zum November fortgesetzt und dann aufgegeben. Bei abnehmender Helligkeit des Planeten erlaubte die kleine Öffnung des Fernrohrs an unserem Mauerquadranten keine sehr genaue Bestimmung mehr. Weniger genaue Beobachtungen aus dieser Zeit wären für die Verbesserung der Elemente ohnehin von geringem Nutzen gewesen. Erst gegen Ende dieser Sichtbarkeitsperiode wurden die Beobachtungen deshalb mit dem Kreismikrometer wieder aufgenommen. Das unbeständige Wetter ließ zwar nur zwei Beobachtungen zu, doch sie fielen besonders gut aus. Sie verdienen hier umso eher mitgeteilt zu werden, als es wahrscheinlich die spätesten Beobachtungen dieser Periode waren.


\begin{center}\fontsize{10}{12}\selectfont\setlength{\tabcolsep}{5pt}\renewcommand{\arraystretch}{1.15}\begin{tabular}{lrrr}
1809 & \shortstack{Mittlere Zeit\\in Göttingen} & \shortstack{Scheinbare\\Rektaszension der Vesta} & \shortstack{Scheinbare südliche\\Deklination}\\
\hline
Febr. 8 & $7^{\text{u}}\;8^{\text{m}}\;25^{\text{s}}$ & $9^\circ\;26'\;51''3$ & $2^\circ\;9'\;28''8$\\
16 & $7\;16\;56$ & $12\;25\;20.0$ & $0\;41\;32.4$\\
\end{tabular}\end{center}


\par\indent Die im 107. Stück unserer vorjährigen Anzeigen [1808. Juli 4. S. 302 dieses Bandes] abgedruckten Elemente weichen von diesen letzten Beobachtungen um 4 Minuten in der Länge und um nur wenige Sekunden in der Breite ab. Sie brauchten daher nur mäßig verbessert zu werden, um sie mit den Beobachtungen von 1807, 1808 und 1809 in Übereinstimmung zu bringen. Hier folgen die neuen Elemente, die Professor \textsc{Gauss} kürzlich bestimmt hat.



\Needspace{20\baselineskip}
\SourceLeaf{316}\begin{center}
IV. \emph{Elemente der Vesta.}
\end{center}


\begin{center}\fontsize{10}{12}\selectfont\setlength{\tabcolsep}{5pt}\renewcommand{\arraystretch}{1.15}\begin{tabular}{lr}
\multicolumn{2}{c}{Epochen der mittleren Länge.}\\
\multicolumn{2}{c}{Meridian von Göttingen.}\\
1807 .\;.\;.\;.\;.\;.\;.\;.\;.\; & $168^\circ\;16'\;35''5$\\
1808 .\;.\;.\;.\;.\;.\;.\;.\;.\; & $267\;39\;28.6$\\
1809 .\;.\;.\;.\;.\;.\;.\;.\;.\; & $6\;46\;4.2$\\
1810 .\;.\;.\;.\;.\;.\;.\;.\;.\; & $105\;52\;39.7$\\
1811 .\;.\;.\;.\;.\;.\;.\;.\;\textsuperscript{.}\; & $204\;59\;15.3$\\
\end{tabular}\end{center}

\begin{center}\fontsize{9}{11}\selectfont\setlength{\tabcolsep}{5pt}\renewcommand{\arraystretch}{1.15}\begin{tabular}{lr}

Perihel 1807, fest gegenüber den Fixsternen .\; & $249^\circ\;52'\;23''8$\\
Aufsteigender Knoten 1807, fest gegenüber den Fixsternen  & $103\;13\;11{,}2$\\
Neigung der Bahn .\;.\;.\;.\;.\;.\; & $7\;8\;18{,}8$\\
Tägliche mittlere tropische Bewegung .\;.\; & $977''5221$\\
Tropische Umlaufszeit .\;.\; & 1325 Tage 19 Stunden\\
Exzentrizität .\;.\;.\;.\;.\;.\;.\;.\;.\;.\; & $0{,}0887809$\\
Logarithmus der großen Halbachse .\;.\; & $0{,}3732940$\\
\end{tabular}\end{center}


\par\indent Nach diesen Elementen erwarten wir die zweite Opposition der Vesta am Nachmittag des 1. Januar 1810 bei $100^\circ\;43'$ Länge und $0^\circ\;31'$ südlicher geozentrischer Breite.

\Needspace{18\baselineskip}
\par\indent Vielleicht finden einige unserer Leser auch die Ergebnisse über die wechselnde relative Helligkeit dieses Planeten interessant, die Professor \textsc{Gauss} bei diesen Rechnungen nebenbei ermittelt hat. Als Einheit dient die Helligkeit, die der Planet im Abstand 1 von Erde und Sonne zeigen würde. Am Tag seiner Entdeckung, dem 29. März 1807, als der Planet wie ein Stern 5..6. Größe erschien, war die Helligkeit $= 0,11815$. Sie nahm damals schon ab und war am 26. September 1807, als die Mailänder Astronomen den Planeten zum letzten Mal beobachteten, auf 0.02699 gesunken. Am 22. Juni 1808, als der Planet an unserer Sternwarte erstmals wiedergefunden wurde und wie ein Stern 8. oder 7..8. Größe aussah, war die Helligkeit schon wieder $= 0.05492$. Der Planet wurde immer heller. Im September, als seine Helligkeit bei der Opposition auf 0,09513 gestiegen war, glich er einem Stern 6. oder 6..7. Größe. Danach nahm die Helligkeit wieder ab. Am 16. Februar war sie auf 0,01629 gesunken, obwohl der Planet immer noch gut die Helligkeit eines Sterns 9. Größe hatte. Bei der nächsten Opposition wird sie nur auf 0,06435 steigen. Daher ist der Planet nur mit der Helligkeit eines Sterns 7. Größe zu erwarten.
